% ---------------------------------------------------------------------
%  00-abstract.tex -- abstract text only (wrapped by main.tex)
% ---------------------------------------------------------------------
For more than seven years I have managed complete products as a senior
product manager: hardware, firmware and software, across a large portfolio
with constant rotation of references, improvement requests and competitive
pressure. Backlogs, Kanban, Scrum, roadmaps, Stage-Gate and PLM change
orders each solve part of that job, but none answers the questions I am
asked every week: from which version did this work start, what else is
running on the same product, and which release actually closed it.
Software teams answer exactly these questions for code with
branch-and-merge. The \emph{Horizontal Product Branch} (HPB) methodology
applies branch-and-merge to the product, not to the code: every piece of
work is a typed branch that forks from a released version and ends in a
named release or in a discard with a written reason. This paper defines
HPB through seven rules, a standard notation, an operating guide and
metrics that fall out of the graph. It also presents HPB Markup, an
open-source text language that draws HPB diagrams from a few lines of text
in HTML, React, Next.js, Vite or the command line, and a prototype built
into a product management portal.
